% Archivo compartido: no compilar directamente.
% Use clase_2_estudiante.tex o clase_2_docente.tex como archivo raíz.
\ifdefined\documentoSemanal
\else
\usepackage{tikz}
\usepackage{estilo_apuntes}
\usetikzlibrary{babel}
\encabezado{Semana 1 · Clase 2}{Material de clase}

\begin{document}
\fi
\tituloclase{Representaciones y funciones algebraicas}{Dominio, imagen y formas de representación}{Semana 1 · Clase 2 · 2 horas}

\begin{objetivos}
\begin{itemize}
  \item Relacionar representaciones verbales, tabulares, gráficas y algebraicas.
  \item Identificar las principales clases de funciones algebraicas.
  \item Determinar el mayor conjunto de números reales sobre el cual una ley de asignación define una función real.
  \item Describir el recorrido mediante propiedades algebraicas o gráficas.
\end{itemize}
\end{objetivos}

\begin{resumen}
Una función $f:A\to B$ queda determinada por su dominio, su codominio y su regla de asignación. Puede describirse verbalmente, mediante una tabla, con un gráfico o por una fórmula; cada representación destaca información diferente de la misma correspondencia entrada--salida.
\end{resumen}

\section*{Representaciones de una función}

\begin{definicion}[representaciones]
\begin{itemize}
  \item \textbf{Verbal:} describe la dependencia entre cantidades.
  \item \textbf{Tabular:} organiza pares entrada--salida.
  \item \textbf{Gráfica:} representa los puntos $(x,f(x))$.
  \item \textbf{Algebraica:} expresa la regla mediante una fórmula o por tramos.
\end{itemize}
\end{definicion}

Usaremos una misma situación para comparar las cuatro representaciones. Un depósito contiene inicialmente $1200$ L y pierde agua a una razón constante de $25$ L por minuto.

\begin{ejemplo}[representación verbal]
El volumen depende del tiempo: comienza en $1200$ L y disminuye $25$ L por cada minuto transcurrido, hasta que el depósito queda vacío.
\end{ejemplo}

\begin{ejemplo}[representación tabular]
Algunos pares tiempo--volumen son
\[
\begin{array}{c|ccccc}
t\text{ (min)}&0&12&24&36&48\\ \hline
V\text{ (L)}&1200&900&600&300&0
\end{array}
\]
La tabla permite consultar valores concretos y observar que, cada $12$ minutos, el volumen disminuye $300$ L.
\end{ejemplo}

\begin{ejemplo}[representación gráfica]
Los pares de la tabla pertenecen al segmento que une $(0,1200)$ con $(48,0)$:
\begin{center}
\begin{tikzpicture}[x=0.09cm,y=0.003cm,every node/.style={font=\scriptsize}]
  \draw[Gris] (0,0)--(53,0) node[right] {$t$ (min)};
  \draw[Gris] (0,0)--(0,1320) node[above] {$V$ (L)};
  \foreach \x in {0,12,24,36,48}{
    \draw[Gris] (\x,25)--(\x,-25) node[below] {\x};
  }
  \foreach \y in {300,600,900,1200}{
    \draw[Gris] (0.7,\y)--(-0.7,\y) node[left] {\y};
  }
  \draw[Azul,very thick] (0,1200)--(48,0);
  \foreach \x/\y in {0/1200,12/900,24/600,36/300,48/0}{
    \fill[Azul] (\x,\y) circle (2.2pt);
  }
\end{tikzpicture}
\end{center}
La gráfica permite apreciar de inmediato la disminución constante, el valor inicial y el instante en que se vacía el depósito.
\end{ejemplo}

\begin{ejemplo}[representación algebraica]
La regla que reúne toda la información es
\[
V:[0,48]\to[0,1200],\qquad V(t)=1200-25t.
\]
La fórmula permite calcular cualquier valor, resolver ecuaciones y analizar la función simbólicamente. El intervalo $[0,48]$ proviene del contexto físico.
\end{ejemplo}

\begin{nota}
Una representación puede destacar ciertos aspectos y ocultar otros. Para especificar completamente una función deben conocerse su dominio, su codominio y su regla de asignación. Una fórmula aislada no determina por sí sola el dominio impuesto por una situación real.
\end{nota}

Los \textbf{ceros} de una función son las entradas $x$ para las cuales $f(x)=0$; gráficamente son las intersecciones con el eje horizontal. Conocerlos permite localizar cuándo una cantidad se anula, factorizar expresiones, construir tablas de signos y separar intervalos donde la función es positiva o negativa. En el modelo anterior,
\[
V(t)=0\iff1200-25t=0\iff t=48,
\]
por lo que el cero indica el instante exacto en que el depósito queda vacío. Si $0$ pertenece al dominio, $f(0)$ da la intersección con el eje vertical y suele representar el valor inicial.

\section*{Conjunto dominio de una ley de asignación}

\begin{proposicion}[para determinar el conjunto dominio]
Si solo se proporciona una ley de asignación real $F(x)$, sin declarar una función, se denomina \textbf{dominio natural de la ley} al mayor conjunto
\[
D_F=\{x\in\mathbb{R}:F(x)\text{ está definida como número real}\}.
\]
Para construir $D_F$ se imponen simultáneamente las condiciones de cada subexpresión:
\begin{itemize}
  \item \textbf{Denominadores distintos de cero.} Por ejemplo,
  \[
  f(x)=\frac{x-3}{x^2+x-6}
  \quad\Longrightarrow\quad x\neq-3,\;x\neq2.
  \]
  \item \textbf{Radicandos de raíces pares no negativos.} Por ejemplo,
  \[
  g(x)=\sqrt{5-2x}
  \quad\Longrightarrow\quad 5-2x\geq0,
  \quad x\leq\frac52.
  \]
  \item \textbf{Radicandos de raíces pares en denominadores estrictamente positivos.} Por ejemplo,
  \[
  h(x)=\frac{1}{\sqrt{x^2-9}}
  \quad\Longrightarrow\quad x^2-9>0,
  \quad x<-3\ \text{o}\ x>3.
  \]
  \item \textbf{Raíces impares sin restricción sobre el radicando.} Por ejemplo, $\sqrt[3]{x-4}$ está definida para todo $x\in\mathbb{R}$.
  \item \textbf{Restricciones anidadas desde el interior hacia el exterior.} Por ejemplo,
  \[
  p(x)=\sqrt{\frac{x-1}{x+2}}
  \quad\Longrightarrow\quad x\neq-2
  \quad\text{y}\quad\frac{x-1}{x+2}\geq0.
  \]
\end{itemize}
Sobre $D_F$, la ley de asignación define una función $f:D_F\to\mathbb{R}$ mediante $f(x)=F(x)$. El conjunto $D_F$ es la intersección de todas las condiciones. En un modelo, el contexto puede imponer un dominio más pequeño.
\end{proposicion}

\begin{ejemplo}[dominio con varias restricciones]
Determine el mayor conjunto $D\subseteq\mathbb{R}$ para que la ley de asignación
\[
g(x)=\sqrt{\frac{x-1}{x+2}}
\]
defina una función real $g:D\to\mathbb{R}$.
\end{ejemplo}
\ifmostrarSoluciones
\begin{solucion}
Se exige $(x-1)/(x+2)\geq0$ y $x\neq-2$. El análisis de signos en $-2$ y $1$ da
\[
\operatorname{Dom}(g)=(-\infty,-2)\cup[1,\infty).
\]
\end{solucion}
\else\vspace{15mm}\fi

\section*{Funciones elementales}

\begin{definicion}[funciones lineal, afín y cuadrática]
Dados $a,b,c\in\mathbb{R}$, se definen sobre $\mathbb{R}$ las funciones
\[
\ell(x)=ax,\qquad A(x)=ax+b.
\]
La primera se llama \textbf{lineal} y la segunda, \textbf{afín}. Si $a\neq0$, la función
\[
q:\mathbb{R}\to\mathbb{R},\qquad q(x)=ax^2+bx+c,
\]
se llama \textbf{cuadrática}.
\end{definicion}

\begin{ejemplo}[imagen de una cuadrática]
Sea $f:\mathbb{R}\to\mathbb{R}$, $f(x)=x^2-4x+7$. Determine su imagen.
\end{ejemplo}
\ifmostrarSoluciones
\begin{solucion}
Al completar el cuadrado,
\[
f(x)=(x-2)^2+3\geq3.
\]
El mínimo $3$ se alcanza en $x=2$; por tanto, $\operatorname{Im}(f)=[3,\infty)$.
\end{solucion}
\else\vspace{15mm}\fi

\begin{definicion}[función polinómica]
Sean $n\in\mathbb{N}^{*}$ y $a_0,\ldots,a_n\in\mathbb{R}$, con $a_n\neq0$. La función
\[
p:\mathbb{R}\to\mathbb{R},\qquad
p(x)=\sum_{k=0}^{n}a_kx^k
=a_nx^n+\cdots+a_1x+a_0
\]
se llama \textbf{polinómica de grado $n$}.
\end{definicion}

\begin{definicion}[función racional]
Sean $p,q:\mathbb{R}\to\mathbb{R}$ funciones polinómicas, con $q$ no idénticamente nula, y sea
\[
D=\{x\in\mathbb{R}:q(x)\neq0\}.
\]
La función
\[
r:D\to\mathbb{R},\qquad r(x)=\frac{p(x)}{q(x)},
\]
se llama \textbf{racional}.
\end{definicion}

\begin{definicion}[función radical]
Sea $n\geq2$ un entero y defina
\[
D_n=
\begin{cases}
[0,\infty),& n\text{ par},\\
\mathbb{R},& n\text{ impar}.
\end{cases}
\]
La función radical de índice $n$ es
\[
\rho_n:D_n\to\mathbb{R},\qquad \rho_n(x)=\sqrt[n]{x},
\]
donde $\bigl(\sqrt[n]{x}\bigr)^n=x$ y, cuando $n$ es par, $\sqrt[n]{x}\geq0$.
\end{definicion}

\begin{definicion}[función algebraica]
Una función \textbf{algebraica} es aquella que puede expresarse mediante un número finito de operaciones algebraicas entre polinomios, incluyendo sumas, restas, productos, cocientes y extracción de raíces.
\end{definicion}

\begin{ejercicios}[para la clase]
\begin{enumerate}
  \item Para la ley $V(t)=1200-25t$, determine el mayor conjunto dominio real y su imagen; luego establezca el dominio y la imagen impuestos por el contexto físico.
  \item La tabla $\{(-2,4),(-1,1),(0,0),(1,1),(2,4)\}$ sugiere una regla. Proponga su fórmula y extensión natural.
  \item Determine el mayor conjunto $D_f\subseteq\mathbb{R}$ para que $f(x)=\dfrac{x-3}{x^2+x-6}$ defina una función real $f:D_f\to\mathbb{R}$.
  \item Determine el mayor conjunto $D_g\subseteq\mathbb{R}$ para que $g(x)=\sqrt{5-2x}$ defina una función real $g:D_g\to\mathbb{R}$.
  \item Determine el mayor conjunto $D_h\subseteq\mathbb{R}$ para que $h(x)=\dfrac{\sqrt{x+1}}{x-4}$ defina una función real $h:D_h\to\mathbb{R}$.
  \item Determine el mayor conjunto $D_p\subseteq\mathbb{R}$ para que $p(x)=\sqrt{\dfrac{x+3}{x-1}}$ defina una función real $p:D_p\to\mathbb{R}$.
  \item Determine el mayor conjunto $D_q\subseteq\mathbb{R}$ para que $q(x)=\dfrac{1}{\sqrt{x^2-9}}$ defina una función real $q:D_q\to\mathbb{R}$.
  \item Sea $u:\mathbb{R}\to\mathbb{R}$, $u(x)=(x+3)^2-5$. Determine su imagen.
  \item Sea $v:\mathbb{R}\to\mathbb{R}$, $v(x)=4-2x^2$. Determine su imagen.
  \item Sea $w:\mathbb{R}\setminus\{-1\}\to\mathbb{R}$, $w(x)=\dfrac{2}{x+1}-3$. Determine su imagen.
  \item Sea $R:\mathbb{R}\setminus\{-1\}\to\mathbb{R}$, $R(x)=\dfrac{x}{x+1}$. Determine su imagen despejando $x$ en términos de $y$.
\end{enumerate}
\end{ejercicios}

\ifmostrarSoluciones
\begin{solucion}[de los ejercicios]
\begin{enumerate}
  \item El mayor conjunto dominio real de la ley y la imagen de la función así definida son $\mathbb{R}$. En el modelo físico, el dominio y la imagen son $[0,48]$ y $[0,1200]$.
  \item $f(x)=x^2$, con dominio $\mathbb{R}$ e imagen $[0,\infty)$.
  \item $\mathbb{R}\setminus\{-3,2\}$.
  \item $(-\infty,5/2]$.
  \item $[-1,4)\cup(4,\infty)$.
  \item $(-\infty,-3]\cup(1,\infty)$.
  \item $(-\infty,-3)\cup(3,\infty)$.
  \item $[-5,\infty)$.
  \item $(-\infty,4]$.
  \item $\mathbb{R}\setminus\{-3\}$.
  \item De $y=x/(x+1)$ se obtiene $x=y/(1-y)$; la imagen es $\mathbb{R}\setminus\{1\}$.
\end{enumerate}
\end{solucion}
\fi

\fuentes{\textbf{Para ampliar:} \emph{Cálculo en una variable: notas de clase 2018-B}, Resumen No. 1, pp. 1--6; J. Stewart, \emph{Cálculo de una variable: trascendentes tempranas}, §§1.1--1.2; y G. Strang--E. Herman, \emph{Cálculo, volumen 1}, OpenStax, §§1.1--1.2.}
\ifdefined\documentoSemanal
\else
\end{document}
\fi
